% !TEX root = ../../proposal.tex

\section{Research Hypotheses}
\label{sec:research_objectives_research_hypothesis}

The main hypotheses of this proposal are:
\begin{itemize}
    \item \textbf{H1 (The Composability Gap):}
    Logical composition via score addition produces a measurable composability gap relative to semantic composition, emerging early in inference and following a principled ordering: smallest for co-occurring pairs, small-to-moderate for approximately disentangled (e.g style-content pairs) unseen pairs, large for entangled pairs, and largest for semantically competing pairs.
    The largest-gap outcomes should be unreachable via natural-language prompting alone, confirming that logical composition steers generation along paths no single prompt can reproduce.

    \item \textbf{H2 (The Representation Hypothesis):}
    The composability gap is primarily a representation problem.
    Holding the compositional method fixed, a model that keeps concept representations explicitly separate should produce better compositions than one that does not.

    \item \textbf{H3 (The Clinical Independence Sufficiency Hypothesis):}
    Approximate clinical independence between chest X-ray pathologies is sufficient for logical composition to generate plausible co-morbid outputs. This includes out-of-distribution pairings unseen during training, provided the model has learned representations in which pathology-specific information remains sufficiently distinct.

\end{itemize}
